% Part: proof-theory
% Chapter: normalization
% Section: introduction

\documentclass[../../../include/open-logic-section]{subfiles}

\begin{document}

\olfileid{pt}{nor}{int}

\olsection{تمهيد في التطبيع}

إذا أمكننا أن~!!{prove} الشرطية~$!A \lif !B$، جاز أن نستعملها في
!!a{proof} لـ~$!B$ بإعمال~\Elim{\lif}. ولا يتم ذلك إلا مع
!!a{proof}~$\delta_1$ لـ~$!A$. والغالب أن يكون آخر~!!{proof} الموجود
لـ~$!A \lif !B$ هو~\Intro{\lif}، وهي تنقل~!!a{proof}~$\delta_2$
لـ~$!B$ من الفرض المفتوح~$!A$ إلى برهان لـ~$!A \lif !B$. فإن كان
الأمر كذلك، صار~!!{proof} النهائي على الصورة:
\[
\AxiomC{$\Discharge{!A}{x}$}
\RightLabel{$\delta_2$}
\DeduceC{$!B$}
\DischargeRule{\Intro{\lif}}{x}
\UnaryInfC{$!A \lif !B$}
\AxiomC{}
\RightLabel{$\delta_1$}
\DeduceC{$!A$}
\RightLabel{\Elim{\lif}}
\BinaryInfC{$!B$}
\DisplayProof
\]
وهذا التدبير نافع من جهة إعادة استعمال ما حصل بالفعل، أعني~!!a{proof}
لـ~$!A \lif !B$، لكنه يجعل~!!{proof} نفسه مشتملًا على فضول؛ إذ ندخل
$!A \lif !B$ ثم نحذفها من فورنا. فلا بد أن نجد طريقًا أقرب كي~!!{prove}~$!B$.

وهذا الطريق موجود أبدًا. فالالتفاف من النوع المتقدم---أي~!!a{introduction}
تتعقبها من فورها قاعدة~!!a{elimination}---يجوز حذفه دائمًا من~!!{proof}s
الاستنباط الطبيعي. وتسمى العملية \emph{التطبيع}، وما ينتج منها~!!{proof}
\emph{سوي} (أو~!!a{proof} \emph{في صورة سوية}).

وبرهان هذه النتيجة شبيه بمبرهنة حذف القطع في حساب المتتاليات: فيه شيء من
التعقيد وكثرة الحالات، وإقامته للنسق الكلاسيكي~\Log{N1c} أشق من إقامته
للنسق الحدسي~\Log{N1i}. فكما يدل حذف القطع على أنه لا مزيد مما~!!{prove}
بإضافة قاعدة~\Cut{} إلى حساب المتتاليات، يدل التطبيع على أن الالتفافات لا
تزيد ما يمكن إثباته على ما يثبت بغيرها. وبينهما شبه آخر أيضًا: فقد يطول
!!{proof} السوي كثيرًا عن أقصر~!!{proof} ذي التفافات للنتيجة نفسها، كما
قد تقصر~!!{proof}s حساب المتتاليات ذات القطوع كثيرًا عن أقصر~!!{proof}
خال من القطع. ثم إن خاصية~!!{formula} الفرعية تقضي بأننا نستطيع دائمًا
أن~!!{prove} النتيجة بما لا يجاوز~!!{formula}s الفرعية لصيغها أو لفروض
!!{proof} المفتوحة؛ وهي لازمة عن تطبيع الاستنباط الطبيعي كما تلزم عن
حذف القطع في حساب المتتاليات.

\end{document}
